\documentclass[11pt]{article}

\usepackage{amsmath, amssymb}
\usepackage{geometry}
\usepackage{bm}
\usepackage{physics}
\usepackage{graphicx}

\geometry{margin=1in}

\title{Governing Equations for Storm Surge Modeling\\
Depth-Averaged Nonlinear Shallow Water Equations}
\author{}
\date{}

\begin{document}

\maketitle

\section{Introduction}

Storm surge in coastal regions is commonly modeled using the depth-averaged nonlinear shallow water equations (2D SWE). These equations describe long-wave, barotropic flow where horizontal length scales are much larger than the vertical scale:

\[
\frac{H}{L} \ll 1.
\]

Under this assumption, vertical acceleration is negligible compared to gravity, and the hydrostatic approximation is valid.

\section{Variables and Definitions}

Let

\begin{itemize}
    \item $\eta(x,y,t)$ = free surface elevation relative to mean sea level,
    \item $h(x,y)$ = bathymetric depth (positive downward),
    \item $H = h + \eta$ = total water depth,
    \item $\bm{u} = (u,v)$ = depth-averaged horizontal velocity,
    \item $\bm{q} = H\bm{u}$ = volume flux (transport per unit width),
    \item $\rho_w$ = water density (assumed constant),
    \item $g$ = gravitational acceleration.
\end{itemize}

\section{Mass Conservation}

Depth-integrated continuity yields:

\begin{equation}
\frac{\partial H}{\partial t}
+
\nabla \cdot (H \bm{u})
=
0.
\end{equation}

This equation states that changes in water depth are governed by horizontal transport divergence.

\section{Momentum Conservation}

The depth-averaged horizontal momentum equations are:

\begin{equation}
\frac{\partial (H\bm{u})}{\partial t}
+
\nabla \cdot
\left(
H\bm{u}\otimes\bm{u}
+
\frac{1}{2} g H^2 \bm{I}
\right)
=
\bm{S},
\end{equation}

where $\bm{I}$ is the identity tensor and $\bm{S}$ contains physical source terms relevant to storm surge.

\section{Physical Source Terms}

The source term includes bathymetry, atmospheric forcing, Coriolis effects, and friction:

\begin{align}
\bm{S}
&=
- g H \nabla h
- \frac{H}{\rho_w} \nabla p_{atm}
+ \frac{\bm{\tau}_w}{\rho_w}
+ f \hat{\bm{k}} \times (H\bm{u})
- C_f |\bm{u}| (H\bm{u})
+ \nabla \cdot (\nu H \nabla \bm{u}).
\end{align}

Each term represents:

\subsection*{Bathymetric Slope}
\[
- g H \nabla h
\]
Ensures correct balance at rest (well-balanced property).

\subsection*{Atmospheric Pressure Gradient}
\[
- \frac{H}{\rho_w} \nabla p_{atm}
\]
Represents inverse barometer effect due to low pressure.

\subsection*{Wind Stress}
\[
\bm{\tau}_w = \rho_{air} C_D |\bm{U}_{10}| \bm{U}_{10}
\]
Surface wind forcing that drives storm surge.

\subsection*{Coriolis Force}
\[
f \hat{\bm{k}} \times (H\bm{u}),
\qquad
f = 2 \Omega \sin\phi
\]
Important for large-scale coastal basins.

\subsection*{Bottom Friction}
Quadratic drag:
\[
- C_f |\bm{u}| (H\bm{u})
\]
or Manning formulation:
\[
- g n^2 \frac{|\bm{u}|}{H^{1/3}} (H\bm{u}).
\]

\subsection*{Horizontal Viscosity}
\[
\nabla \cdot (\nu H \nabla \bm{u})
\]
Represents unresolved turbulence.

\section{Final Governing System}

The complete depth-averaged nonlinear shallow water system for storm surge is:

\begin{equation}
\boxed{
\begin{aligned}
\frac{\partial H}{\partial t}
+
\nabla \cdot (H\bm{u})
&= 0, \\
\frac{\partial (H\bm{u})}{\partial t}
+
\nabla \cdot
\left(
H\bm{u}\otimes\bm{u}
+
\frac{1}{2} g H^2 \bm{I}
\right)
&=
- g H \nabla h
- \frac{H}{\rho_w} \nabla p_{atm}
+ \frac{\bm{\tau}_w}{\rho_w}
+ f \hat{\bm{k}} \times (H\bm{u})
- C_f |\bm{u}| (H\bm{u}).
\end{aligned}
}
\end{equation}

\section{Wave Speed and CFL Condition}

The characteristic gravity wave speed is:

\[
c = \sqrt{gH}.
\]

Time stepping must satisfy a CFL condition based on this speed.

\section{Physical Interpretation}

Storm surge arises from:

\begin{itemize}
    \item Wind stress pushing water toward the coast,
    \item Low atmospheric pressure elevating the sea surface,
    \item Coastal boundaries trapping water,
    \item Nonlinear shallow water dynamics amplifying long waves.
\end{itemize}

The shallow water equations capture these effects through conservative mass and momentum balances.

\end{document}